\documentclass[10pt,a4paper]{amsart}
\usepackage[margin=19mm]{geometry}
\usepackage{amsmath,amssymb,mathtools}
\usepackage[scr=rsfs,cal=euler]{mathalfa}
\usepackage{xfrac}
\usepackage[unicode,hidelinks,bookmarks=false]{hyperref}
\newcommand{\ie}{i.e.\ }
\newcommand{\cf}{cf.\ }
\newcommand{\resp}{respectively\ }
\newcommand{\Subsection}{\subsection}
\providecommand{\nobreakdash}{\nobreak\hbox{-}}
\setlength{\emergencystretch}{2em}
\multlinegap0pt
\usepackage{bm}
\usepackage{bbm}
\usepackage{dsfont}
\usepackage{xspace}
\usepackage{cancel}
\usepackage{mathtools}
\usepackage{amsthm}
\makeatletter
\@ifundefined{proposition}{\newtheorem{proposition}{Proposition}}{}
\@ifundefined{theorem}{\newtheorem{theorem}[proposition]{Theorem}}{}
\makeatother
\DeclareMathOperator{\sing}{Sing}
\newcommand{\C}{\mathbb{C}}
\newcommand{\F}{\mathcal{F}}
\newcommand{\W}{\mathbf{W}}
\newcommand{\oo}{\mathcal{O}}
\title[A lower bound for the Milnor number of vector fields]{A lower bound for the Milnor number of vector fields}
\author{Maur\'icio Corr\^ea, Gilcione Nonato Costa, Alejandra Salamanca Russi}
\date{Source: arXiv:2602.10047v2 (CC BY 4.0)}
\begin{document}
\maketitle

\noindent\emph{Source and licence.} A lower bound for the Milnor number of vector fields. Version arXiv:2602.10047v2 (CC BY 4.0), distributed under the Creative Commons Attribution 4.0 International licence, as recorded in the arXiv licence metadata for this exact version and shown on the abstract page https://arxiv.org/abs/2602.10047v2. Full rights notice: licenses/arxiv-2602.10047-CC-BY-4.0.txt.\par\smallskip

\noindent\textit{Editorial scope.} Counted unit: the Theorem whose source label is \texttt{thm:polynomial-approximation} in the article (source0.tex lines 619--627, source label \texttt{thm:polynomial-approximation}), delivered together with the article's own complete original proof of it (source0.tex lines 629--717, one \texttt{proof} environment of 89 lines). No mathematics is written, repaired or re-derived: statement and proof are the article's own text line for line. Required context retained uncounted: none. Every definition, display and label of those lines is retained unchanged. Omitted from this package and not redistributed: 1--618, 628--628, 718--1337. Nothing inside a retained excerpt is omitted or altered: every retained body is delivered exactly as the article's lines are written, so no omission receipt of any kind accompanies this package. Citations: the retained excerpts carry no citation command at all, so no bibliography entry of the article is transcribed and no reference is redistributed. Reference adaptation: none was needed and none was made; every reference inside the delivered unit resolves inside it, so no unresolved reference or citation is produced. Printed slips kept literal: no printed word, symbol, hypothesis or number of the counted statement or of its proof was altered. Layout and class: the article's own class is \texttt{amsart} with packages including fontenc, inputenc, stix2, microtype, geometry, graphicx, mathtools, amsthm, mathrsfs; the wrapper is the lane's pinned \texttt{amsart} editorial wrapper at 10pt on A4, which restates the theorem-like environments the retained text needs and adds the article's own macro definitions that the retained text needs (\texttt{\textbackslash C}, \texttt{\textbackslash F}, \texttt{\textbackslash W}, \texttt{\textbackslash oo}, \texttt{\textbackslash sing}), with extra wrapper packages (bm, bbm, dsfont, xspace, cancel, mathtools). The delivered numbering is the unit's own. Assets: no retained span contains a figure environment, a graphics-inclusion command, a PDF-page inclusion, a table or a TikZ picture; no image file is copied into this package, the package declares no asset dependency and no image file is redistributed with it. Licence: version arXiv:2602.10047v2 of this article is distributed under the Creative Commons Attribution 4.0 International licence, as recorded in the arXiv licence metadata for this exact version and shown on the abstract page https://arxiv.org/abs/2602.10047v2. A full rights notice is delivered at licenses/arxiv-2602.10047-CC-BY-4.0.txt. Characters: every retained excerpt is pure ASCII, so the delivered engine, XeLaTeX with its default Latin Modern text font, reports no missing character.
\begin{theorem}\label{thm:polynomial-approximation}
Let $\F$ be a one-dimensional holomorphic foliation on $\C^n$ with
$\W=Z(f_1,\ldots,f_d)\subset\sing(\F)$ a smooth complete intersection on $\C^n$ and each $f_i\in \oo_{\C^n}$.
Then there are polynomial approximations $\{\F_\kappa\}$ of $\F$ and analytic subsets $W_\kappa\subset\sing(\F_\kappa)$ such that:
\begin{enumerate}
\item[(a)] $\W_{\kappa}$ is a smooth algebraic local complete intersection for all $\kappa$;
\item[(b)] $\displaystyle\lim_{\kappa\to \infty}\W_{\kappa}=\W$ on each precompact $\epsilon$-neighbourhood of $W$.
\end{enumerate}
\end{theorem}

\begin{proof}
Let $X:=\{X_\alpha,U_\alpha\}$ be a vector field defining $\F$ in local coordinates $\{U_\alpha,\varphi_\alpha\}$, and write
\begin{equation}\label{reference_field}
X(z):=\sum_{i=1}^n X_i(z)\frac{\partial}{\partial z_i}=(X_1(z),\ldots,X_n(z)).
\end{equation}
Let $f_i\in \oo_{\C^n}$ be germs of holomorphic functions such that
$\W:=Z(f_1,\ldots,f_d)$ is a smooth complete intersection subvariety.
After reordering the coordinates if necessary, we may assume that a $d\times d$ minor of the Jacobian matrix $Jf$ of $f$
is invertible on a neighbourhood $U$ of some point of $W$.
Hence the map $\varphi:\C^n\to\C^n$ given by
\[
z\longmapsto \varphi(z)=(f_1(z),\ldots,f_d(z),z_{d+1},\ldots,z_n)=w
\]
is a local biholomorphism.
The push-forward $Y=\varphi_*X$ is a holomorphic vector field such that
$\W_0:=\varphi(W)=Z(w_1,\ldots,w_d)\subset\sing(Y)$ is the common zero locus of the coordinate functions $w_i$ for $i=1,\ldots,d$, and
\begin{equation}\label{eq:push-forward-field}
Y(w):=\varphi_*X(u)=\sum_{i=1}^n Y_i(w)\frac{\partial}{\partial w_i}.
\end{equation}

Fix a relatively compact open $\epsilon/2$-neighbourhood $V\subset\C^n$ such that the set of embedded points $\mathcal{A}_{\W}$
is contained in $V$.
After shrinking $V$ if necessary, we may assume that $\overline V$ is contained in a polydisc centred at the origin on which all coefficient functions $X_i$
in \eqref{reference_field} admit convergent power series expansions.
Let $X_\kappa$ be the holomorphic vector field obtained by truncating the power series of each coefficient $X_i$ to total degree $\leq \kappa$.
Then $X_\kappa$ is polynomial and $X_\kappa\to X$ uniformly on $\overline V$ as $\kappa\to\infty$.
For a fixed sufficiently small $\epsilon>0$, choose $\kappa$ so large that $X_\kappa$ is $\epsilon/2$-close to $X$ on $\overline V$, and retain this value of $\kappa$.

Let $Y=\varphi_*X$ be the vector field inducing the push-forward one-dimensional foliation $\mathcal{G}$, where
$\varphi(z)=(f_1(z),\ldots,f_d(z),z_{d+1},\ldots,z_n)$.
Keeping this notation, denote by $\mathcal{G}_\kappa$ the foliation induced by the polynomial vector field $Y_\kappa$
associated with $X_\kappa$ via the corresponding truncation $\varphi_\kappa$.
Since $\W_0=\varphi(\W)$ is a smooth complete intersection in $\C^n$, its projective closure, obtained by homogenisation, is a smooth complete intersection in $\mathbb P^n$.

To construct the required polynomial deformation, let $Y_\kappa$ be as in \eqref{eq:push-forward-field} and set
\begin{equation}
Y_{\kappa,t}=Y_\kappa+t\,\widetilde{Y},
\end{equation}
where
\[
\widetilde{Y}=Q_1\frac{\partial}{\partial w_1}+\cdots+Q_n\frac{\partial}{\partial w_n},
\]
and
\[
Q_j(w)=
\begin{cases}
\displaystyle\sum_{|I|=q_j} a_{I,j}\, w_1^{i_1}\cdots w_d^{i_d} & \text{if } 1\leq j\leq d,\\[6pt]
\displaystyle\sum_{|I|=q_j} R_{I,j}(w)\, w_1^{i_1}\cdots w_d^{i_d} & \text{if } d+1\leq j\leq n,
\end{cases}
\]
where $a_{I,j}\in\C$ and each $R_{I,j}$ is an affine linear function.
Impose
\[
q_1=\cdots=q_d=q_{d+1}+1=\cdots=q_n+1=\ell+1,
\]
and choose the coefficients so that $\deg(Q_j)\leq \deg(Y_\kappa)$ and the hyperplane at infinity $H^\infty$ is invariant under $Y_{\kappa,t}$.
After an arbitrarily small perturbation of the coefficients of $\widetilde Y$, we may assume that $Y_{\kappa,t}$ is special along $\W_0$ and that $\mathcal G_{\kappa,t}$ has no embedded points associated with $\W_0$ for $t\neq0$. The deformation then satisfies:
\begin{enumerate}
\item[(i)] $\mathcal{G}_{\kappa,0}=\mathcal{G}_\kappa$;
\item[(ii)] $\deg(\mathcal{G}_{\kappa,t})=\deg(\mathcal{G}_\kappa)$;
\item[(iii)] If $\ell=0$, then $\sing(\mathcal{G}_{\kappa,t})=\{p_1^t,\ldots,p_{s_t}^t\}$ and $W_0$ is $\mathcal{G}_{\kappa,t}$-invariant for every $t\in D_\epsilon\setminus\{0\}$;
\item[(iv)] If $\ell>0$, then $\sing(\mathcal{G}_{\kappa,t})=W_0\cup\{p_1^t,\ldots,p_{s_t}^t\}$ is a disjoint union for every $t\in D_\epsilon\setminus\{0\}$, and $\mathcal{G}_{\kappa,t}$ is special along $W_0$ with $\ell=m_E(\mathcal{G}_{\kappa,t},W_0)$.
\end{enumerate}

The polynomial vector field $Y_{\kappa,t}$ defines a holomorphic foliation on all of $\C^n$. Properties (i) and (ii) follow from the definition; properties (iii) and (iv) follow from a suitable choice of the coefficients $a_{I,j}$ and $R_{I,j}$.

After reducing $\epsilon$ if necessary, choose a polynomial vector field
\[
\alpha=\sum_{i=1}^n \alpha_i(w)\frac{\partial}{\partial w_i},
\qquad
\alpha_i\in \C[w_1,\ldots,w_n],
\qquad
\deg(\alpha)\leq \deg(Y_\kappa),
\]
such that the deformation
\begin{equation}\label{eq:principal-deformation}
\widetilde{Y}_t
=
Y_{\kappa,t}+t\sum_{i=1}^n \alpha_i \frac{\partial}{\partial w_i}
\end{equation}
has only isolated singularities for every $t\ne 0$ on the chosen relatively compact $\epsilon/2$-neighbourhood.
Equivalently, the foliations induced by $\widetilde{Y}_t$ have only isolated singularities on that neighbourhood whenever $t\ne 0$.
For every sufficiently small $t\neq0$, a neighbourhood of each embedded point of $Y$ contains singular points of $\widetilde Y_t$, and their total multiplicity equals the number of embedded points associated with $W_0$. Hence the singularities of $\widetilde Y_t$ converging to $W_0$ form a finite set of constant total multiplicity for $0<|t|<\epsilon/2$.
To prove minimality, let $\{Z_t\}_{t\in D'_\epsilon}$ be any other generic holomorphic deformation of $Y$ whose singularities in the same relatively compact $\epsilon/2$-neighbourhood $V$ are isolated.
Fix a norm $\|\cdot\|$ on $\oo_{\C^n}(V)$.
Since $\|Z_t-Y\|\leq \epsilon/2$ and $\|\widetilde{Y}_t-Y\|\leq \epsilon/2$, it follows that
$\|\widetilde{Y}_t-Z_t\|<\epsilon$ on $V$ for $\epsilon$ sufficiently small.
Hence the number of points of $\sing(Z_t)$, counted with multiplicity, is at least the number of points in $\mathcal A_{\W}$. By continuity, the corresponding isolated singularities converge to $\W_0$.
\end{proof}

\end{document}
